\documentclass[11pt, a4paper]{article}

\usepackage[T1]{fontenc}
\usepackage[utf8]{inputenc}
\usepackage{tgpagella}
\usepackage[spanish, es-noshorthands]{babel}
\usepackage{amsmath, amssymb, bm}
\usepackage{booktabs}
\usepackage[table, xcdraw]{xcolor}
\usepackage[most]{tcolorbox}
\usepackage[american]{circuitikz}
\usepackage[margin=2.5cm]{geometry}
\usepackage{fancyhdr}

\pagestyle{fancy}
\fancyhf{}
\setlength{\headheight}{16pt}
\lhead{\textbf{UCB Sede Tarija} | Ing. Mecatrónica}
\chead{}
\rhead{IMT-121 Circuitos Electrónicos I | Soluciones S06}
\lfoot{Prof: M.Sc. Ing. Bernardo Quiroga Turdera}
\rfoot{Página \thepage}
\renewcommand{\headrulewidth}{0.4pt}

\definecolor{ucbblue}{RGB}{0, 51, 102}

\begin{document}

\begin{tcolorbox}[colback=green!5!white, colframe=green!50!black, title=\textbf{Solucionario: Polarización DC del BJT}]
    \centering \large \textbf{Análisis de Redes de Polarización Fija}
\end{tcolorbox}

\section*{Solución al Ejercicio 1: Reconocimiento de Región}
\begin{enumerate}
    \item \textbf{Corte:} Ambas uniones están sin polarización directa. El transistor está bloqueado.
    \item \textbf{Activa Directa:} Unión BE conduciendo, Unión BC en inversa. Es la región de amplificación normal.
    \item \textbf{Saturación:} Ambas uniones polarizadas en directa. Actúa como un interruptor cerrado (corriente limitada por la red externa).
\end{enumerate}

\section*{Solución al Ejercicio 2: Análisis DC Guiado}
\textbf{Circuito:} $V_{CC}=12\,\text{V}$, $R_B=560\,\text{k}\Omega$, $R_C=2.2\,\text{k}\Omega$, $\beta=100$, $V_{BE}=0.7\,\text{V}$.

\textbf{2. Corriente de base ($I_B$):}
\begin{equation*}
    I_B = \frac{V_{CC} - V_{BE}}{R_B} = \frac{12 - 0.7}{560 \times 10^3} \approx 20.18\,\mu\text{A}
\end{equation*}

\textbf{3. Corriente de colector ($I_C$) asumiendo región activa:}
\begin{equation*}
    I_C = \beta I_B = 100 \times 20.18\,\mu\text{A} \approx 2.018\,\text{mA}
\end{equation*}

\textbf{4. Tensión Colector-Emisor ($V_{CE}$):}
\begin{equation*}
    V_{CE} = V_{CC} - I_C R_C = 12 - (2.018 \times 10^{-3})(2.2 \times 10^3) \approx 7.56\,\text{V}
\end{equation*}

\textbf{5. Consistencia:}
$V_{BC} = V_{BE} - V_{CE} = 0.7 - 7.56 = -6.86\,\text{V} < 0$. La suposición de Activa Directa es \textbf{válida}.

\textbf{6. Recta de Carga:}
\begin{equation*}
    V_{CE} = 12 - I_C(2.2\,\text{k}\Omega)
\end{equation*}

\textbf{7. Intersecciones:}
\begin{itemize}
    \item Corte ($I_C = 0$): $V_{CE} = 12\,\text{V}$
    \item Saturación ideal ($V_{CE} = 0$): $I_C = 12 / 2.2\,\text{k}\Omega \approx 5.45\,\text{mA}$
\end{itemize}

\textbf{8. Punto Q:}
$\mathbf{Q \approx (2.02\,\text{mA}, 7.56\,\text{V})}$

\textbf{9. Predicciones:}
\begin{itemize}
    \item Si $R_B \uparrow$: $I_B \downarrow \implies I_C \downarrow \implies V_{CE} \uparrow$ (Q se mueve hacia el corte).
    \item Si $R_C \uparrow$: $I_B$ e $I_C$ (teóricos) no cambian, pero $V_{CE}$ cae (la recta de carga se hace más horizontal). Si $R_C$ aumenta lo suficiente, obligará a la saturación.
\end{itemize}


\section*{Solución al Ejercicio 4: Problema Independiente}
\textbf{Circuito:} $V_{CC}=9\,\text{V}$, $R_B=330\,\text{k}\Omega$, $R_C=1.5\,\text{k}\Omega$, $\beta=120$, $V_{BE}=0.7\,\text{V}$.

\textbf{1. Cálculos de Mallas:}
\begin{equation*}
    I_B = \frac{9 - 0.7}{330 \times 10^3} \approx 25.15\,\mu\text{A}
\end{equation*}
\begin{equation*}
    I_C = 120(25.15\,\mu\text{A}) \approx 3.018\,\text{mA}
\end{equation*}
\begin{equation*}
    V_{CE} = 9 - (3.018\,\text{mA})(1.5\,\text{k}\Omega) \approx 4.47\,\text{V}
\end{equation*}

\textbf{2. Validación:}
$V_{BC} = 0.7 - 4.47 = -3.77\,\text{V} < 0$. La región activa es \textbf{válida}.

\textbf{3 \& 4. Recta de Carga e Intersecciones:}
\begin{equation*}
    V_{CE} = 9 - (1.5\,\text{k}\Omega)I_C
\end{equation*}
Corte: $V_{CE} = 9\,\text{V}$. Saturación ideal: $I_C = 9 / 1.5\,\text{k}\Omega = 6\,\text{mA}$.

\textbf{5. Punto Q:}
$\mathbf{Q \approx (3.02\,\text{mA}, 4.47\,\text{V})}$

\textbf{6. Impacto del cambio de $\beta$ a 60:}
$I_C$ disminuirá a la mitad (aprox. $1.51\,\text{mA}$). Al caer $I_C$, la caída de tensión en $R_C$ disminuye, por lo que $V_{CE}$ aumentará. El punto Q se moverá acercándose al estado de corte.

\end{document}
